% Unicode symbol fallback for the XeLaTeX pandoc pipeline.
%
% XeLaTeX does not fall back automatically: any glyph missing from the main
% (Palatino) or mono (Menlo) font is silently dropped and shows up as an empty
% box / U+FFFD in the PDF. The document text uses a handful of symbols that
% those fonts lack -- circled digits (①②③…), arrows (→ ↔ ⇕…) and geometric
% marks (◎ ○ △ ▼ ✓ …). We route just those code points to Hiragino Sans, a
% system font on macOS that covers all of them, leaving every other character
% (and the body typeface) untouched.
\usepackage{fontspec}
\usepackage{newunicodechar}
\newfontfamily\symbolfallback{Arial Unicode MS}[Scale=MatchLowercase]
\newcommand{\uf}[1]{{\symbolfallback #1}}

% CJK support: Palatino/Menlo have no Japanese glyphs, so without this any
% Japanese text is silently dropped from the PDF. xeCJK routes CJK code points
% to a dedicated font while leaving Latin text in the body typeface. Fonts are
% picked per platform: Hiragino (macOS), Noto CJK (Linux), Yu Mincho (Windows).
\usepackage{xeCJK}
\IfFontExistsTF{Hiragino Mincho ProN}{%
  \setCJKmainfont{Hiragino Mincho ProN}%
  \setCJKsansfont{Hiragino Sans}%
  \setCJKmonofont{Hiragino Sans}%
}{%
  \IfFontExistsTF{Noto Serif CJK JP}{%
    \setCJKmainfont{Noto Serif CJK JP}%
    \setCJKsansfont{Noto Sans CJK JP}%
    \setCJKmonofont{Noto Sans Mono CJK JP}%
  }{%
    \IfFontExistsTF{Yu Mincho}{%
      \setCJKmainfont{Yu Mincho}%
      \setCJKsansfont{Yu Gothic}%
      \setCJKmonofont{Yu Gothic}%
    }{}%
  }%
}

% Enclosed alphanumerics -- circled digits
\newunicodechar{⓪}{\uf{⓪}}
\newunicodechar{①}{\uf{①}}
\newunicodechar{②}{\uf{②}}
\newunicodechar{③}{\uf{③}}
\newunicodechar{④}{\uf{④}}
\newunicodechar{⑤}{\uf{⑤}}
\newunicodechar{⑥}{\uf{⑥}}
\newunicodechar{⑦}{\uf{⑦}}
\newunicodechar{⑧}{\uf{⑧}}
\newunicodechar{⑨}{\uf{⑨}}
\newunicodechar{⑩}{\uf{⑩}}
\newunicodechar{⑪}{\uf{⑪}}
\newunicodechar{⑫}{\uf{⑫}}
\newunicodechar{⑬}{\uf{⑬}}
\newunicodechar{⑭}{\uf{⑭}}
\newunicodechar{⑮}{\uf{⑮}}
\newunicodechar{⑯}{\uf{⑯}}
\newunicodechar{⑰}{\uf{⑰}}
\newunicodechar{⑱}{\uf{⑱}}
\newunicodechar{⑲}{\uf{⑲}}
\newunicodechar{⑳}{\uf{⑳}}

% Arrows
\newunicodechar{←}{\uf{←}}
\newunicodechar{↑}{\uf{↑}}
\newunicodechar{→}{\uf{→}}
\newunicodechar{↓}{\uf{↓}}
\newunicodechar{↔}{\uf{↔}}
\newunicodechar{↕}{\uf{↕}}
\newunicodechar{⇐}{\uf{⇐}}
\newunicodechar{⇒}{\uf{⇒}}
\newunicodechar{⇔}{\uf{⇔}}
\newunicodechar{⇕}{\uf{⇕}}

% Geometric shapes & marks
\newunicodechar{●}{\uf{●}}
\newunicodechar{○}{\uf{○}}
\newunicodechar{◎}{\uf{◎}}
\newunicodechar{◇}{\uf{◇}}
\newunicodechar{◆}{\uf{◆}}
\newunicodechar{△}{\uf{△}}
\newunicodechar{▲}{\uf{▲}}
\newunicodechar{▽}{\uf{▽}}
\newunicodechar{▼}{\uf{▼}}
\newunicodechar{■}{\uf{■}}
\newunicodechar{□}{\uf{□}}
\newunicodechar{▶}{\uf{▶}}
\newunicodechar{◀}{\uf{◀}}
\newunicodechar{✓}{\uf{✓}}
\newunicodechar{✔}{\uf{✔}}
\newunicodechar{✕}{\uf{✕}}
\newunicodechar{✗}{\uf{✗}}
\newunicodechar{・}{\uf{・}}
\newunicodechar{◐}{\uf{◐}}
\newunicodechar{◑}{\uf{◑}}

% Emoji fallbacks: macOS ships no monochrome emoji font and Apple Color Emoji
% does not work under XeLaTeX, so common status emoji are mapped to monochrome
% equivalents instead of being silently dropped.
\newunicodechar{✅}{\uf{✔}}
\newunicodechar{☑}{\uf{✔}}
\newunicodechar{❌}{\uf{✗}}
\newunicodechar{✖}{\uf{✗}}
\newunicodechar{⭕}{\uf{○}}
\newunicodechar{⚠}{\uf{△}}
\newunicodechar{️}{}% U+FE0F variation selector -- render as nothing
